\section{Introduction}\label{sec:introduction}

\subsection{What is a speech signal processing?}\label{subsec:what-is-a-speech-signal?}
Speech is the fundamental means of communication among people in the modern world. It consists of a combination of
sounds that constitute the human voice. The human voice, produced by the articulatory apparatus for the purpose of
linguistic communication, can be mathematically represented as a one-dimensional function of time. \cite{schroeder}

The sound generation process, including the production of human speech, occurs in space through vibrations that induce
changes in air pressure. These fluctuations enable the propagation of sound waves through the medium. Microphones
convert sound waves into electrical signals, such as voltage or current allowing speech signals to be processed and
manipulated in digital form.

Speech processing encompasses the acquisition, storage, transmission, and output of speech signals in digital format.
This field represents a specialized area of research situated at the intersection of digital signal processing and
natural language processing. \cite{priyanka}


%\section{Problem Description}\label{sec:problem-description}
%
%\paragraph{}
%The loss of audio signal fragments in recordings reduces their quality, which makes further processing and utilization complicated.
%Missing segments may contain critical information, particularly in the case of voice recordings, where minor data loss
%critically affect their perception and interpretation.
%
%Machine learning-based speech signal restoration technologies offer efficient approaches for reconstructing and predicting
%lost data contributing to the preservation of the integrity of audio material.
%

%
%\section{Goal}
%\paragraph{}
%The main task is to create software that reconstructs lost speech signal components using machine learning.
%
%
%\section{Methodology}
%\paragraph{}
%To achieve this goal, the following approach will be used:
%\begin{enumerate}
%    \item \textbf{Creating artificially damaged data:}
%    \begin{itemize}
%        \item Creating a dataset that includes audio files with pre-added artificial damage, such as "gaps" or white noise.
%        \item This method allows you to control the damage process and simplifies the creation of a training dataset.
%    \end{itemize}
%    {\color{gray}
%    \item \textbf{Collecting data with real damage:}
%    \begin{itemize}
%
%        \item Using audio files that have been subject to distortion or noise as a result of recording or transmission.
%
%        \item The main difficulty lies in the need for significant computing resources for training models since the
%        \item diversity and volume of data significantly affect the quality of the result.
%    \end{itemize}}
%
%\end{enumerate}
%
%{\color{red}The first approach is expected to be more efficient in terms of resource costs and scalability.
%
%The second approach could possibly be used in the Master's Thesis.}
%
%
%\section{Technology Stack and Resources for the development}
%\begin{itemize}
%    \item \textbf{Python:} The main programming language for implementing the project and integrating all components.
%    \item \textbf{NumPy and SciPy:} Libraries for analyzing and processing audio signals at a low level.
%    \item \textbf{Librosa:} A tool for extracting and pre-processing audio signal characteristics.
%    \item \textbf{Kaldi, TensorFlow, or PyTorch:} Platforms for creating and training machine learning models that
%    can restore lost segments of an audio signal.
%    \item \textbf{OpenSLR:} Speech signals data acquisition.
%\end{itemize}
%
%\section{Result}
%\paragraph{}
%The result of the work will be a software solution that allows to process damaged audio files and restore lost
%segments.
%The source code of the project will be published on the GitHub platform with detailed instructions for installation and use.
%
%\section{Report Structure}
%\begin{enumerate}
%    \item \textbf{Introduction:}
%    \begin{itemize}
%        \item Introduction to the problem, justification of the relevance and purpose of the study.
%    \end{itemize}
%    \item \textbf{Literature Review:}
%    \begin{itemize}
%        \item Description of the characteristics of the audio signal and its nature.
%        \item Review of classical methods of audio signal processing (e.g., filters).
%        \item Analysis of modern approaches based on machine learning.
%    \end{itemize}
%    \item \textbf{Methodology:}
%    \begin{itemize}
%        \item Detailed description of data collection methods, technologies, and algorithms used. Only necessary
%        technical information, not the explanation of the whole code.
%    \end{itemize}
%    \item \textbf{Results:}
%    \begin{itemize}
%        \item Evaluation of the model results: quality of restoration, testing on data with a specific type of damage.
%    \end{itemize}
%    \item \textbf{Conclusion:}
%    \begin{itemize}
%        \item Results of the work, and software quality description.
%    \end{itemize}
%\end{enumerate}
